% ch12: 原文 GLDraft_D.tex 行 2843–3169；§8
% 只译散文（标题、正文、图题、脚注）；公式、\label、\ref、\eqref、\cite、图文件名保持不变。

\section{Hecke 算符与谱流模}\label{sec:chiral}

\subsection{预备知识}

与增广 Neumann 边界条件（enriched Neumann boundary condition）相联系的 BAA 型膜的量子化，已经为我们提供了波函数的一些例子（\ref{eq:ZPsi}），这些波函数定义为具有手征与反手征临界 Kac--Moody 对称性的二维共形场论的配分函数。本节说明如何直接用二维共形场论的语言，计算量子 Hitchin 哈密顿量与 Hecke 算符作用在这类配分函数上的结果。与此同时，我们也将对第 \ref{wtw} 节中为定义 Hecke 算符而援引的数学结果有更好的理解。

我们已经简要讨论过量子 Hitchin 哈密顿量的二维共形场论解释。临界 Kac--Moody 手征代数有一个很大的中心，它由某些局部算符 $S_{\cal P}(z)$ 生成，这些算符与各流的 OPE 非奇异。关联函数在复化规范变换下的变换方式，由各流的 Ward 恒等式描述。“$S_{\cal P}(z)$ 与各流的 OPE 非奇异”这一断言意味着：插入了这类算符的关联函数只有
\begin{equation}
\langle S_{{\cal P}_1}(z_1) \cdots S_{{\cal P}_n}(z_n) \rangle_a
\end{equation}
满足与配分函数相同的变换性质（\ref{eq:cogauge}）。因此它也定义了 $\M(C,G)$ 上的一个半密度。

此外，$S_{\cal P}(z)$ 由 Kac--Moody 流组装而成，而这些流可以换成关于联络的泛函导数。于是我们可以递推地展开
\begin{equation}
\langle S_{{\cal P}_1}(z_1) \cdots S_{{\cal P}_n}(z_n) \rangle_a = {\cal D}^{(a)}_{\P_1}(z_1 )\langle S_{{\cal P}_2}(z_2) \cdots S_{{\cal P}_n}(z_n) \rangle_a
\end{equation}
而最终结果不依赖于施行递推时所作用算符的次序。因此微分算子 ${\cal D}^{(a)}_{\P_1}(z_1 )$ 彼此交换。

尽管这里谈的是某个二维共形场论的关联函数，但这并非必要：只要给定 $\M(C,G)$ 上的一个半密度，它由联络空间上的规范协变泛函 $\Psi(a)$ 表示，则关于 $a$ 的泛函导数就表现得与 Kac--Moody 流一样。微分算子 ${\cal D}^{(a)}_{\P}(z)$ 以规范协变的方式，表示量子 Hitchin 哈密顿量 ${\cal D}_{\P}(z)$ 对半密度 $\Psi$ 的作用。

在 $C$ 中一点 $p$ 的邻域 $U$ 内作计算时，通常有帮助的做法是为丛选取一个在 $U$ 上为零的代表联络。这总是可能的，因为 $\cA_{\bar z}$ 可以借助取复值的规范变换局部地置零。这相当于把丛在 $U$ 上平凡化，正如我们在第 \ref{affgr} 节讨论一般 Hecke 变换时所做的。于是 Kac--Moody 流在 $U$ 上亚纯，并满足标准形式的 Kac--Moody OPE
\begin{equation}
J^a(z) J^b(w) \sim \frac{\kappa_c \delta^{ab}}{(z-w)^2} + \frac{f^{ab}_d J^d(w)}{z-w}.
\end{equation}

回顾 Kac--Moody 代数 Fourier 模的定义
\begin{equation}
J^a_n \equiv \oint_{|z|=\epsilon} \frac{\d z}{2 \pi \i} z^n J^a(z)
\end{equation}
插入这样一个 Fourier 模，表示 $a_{\bar z}$ 的一个支撑在回路 $|z|=\epsilon$ 上的无穷小形变；或者说，它表示丛的一个形变：借 $U'=U\backslash p$ 中的一个无穷小规范变换，修改 $U$ 上的丛与曲面其余部分上的丛之间的粘合。

若圆盘 $|z|<\epsilon$ 内没有其他算符插入（或只有中心算符插入），则插入了非负 Fourier 模的关联函数为零。相应的无穷小规范变换可以延拓到 $U$ 上，并不改变丛。它们表示 $U$ 上原来那个平凡化的改变。若圆盘内有一个一般的算符插入，则非负模的作用非平凡：插入一个一般的局部算符需要选定丛的某种平凡化，而结果依赖于这一选取。

负的 Fourier 模即使在没有其他算符插入时也可以非平凡地作用，它们表示能够改变丛的无穷小规范变换。负模的反复作用生成 Kac--Moody 代数真空模在 $z=0$ 处的像。算符 $S_{\cal P}(0)$ 以及手征代数中其他中心元，按算符–态对应，就是真空模中被各流的所有非负 Fourier 模消灭的那些向量。例如，Sugawara 向量为
\be
|S_2\rangle \equiv \Tr \,J_{-1} J_{-1} |0 \rangle,
\ee 对其他中心元也有类似的公式。

\subsection{作为中心顶点算子的 Hecke 算符}\label{cvo}

Hecke 积分算符也可以用二维手征代数的方法来分析。我们希望把 Hecke 算符提升为作用在规范协变泛函上的算符，并用与 Kac--Moody 流 OPE 平凡的局部算符的插入，赋予它们二维手征代数的解释。这样的表述立刻保证 Hecke 算符与量子 Hitchin 哈密顿量交换，并且彼此交换。

我们可以逐字照搬第 \ref{affgr} 节中 Hecke 修正的定义。首先，在点 $p$ 的一个小邻域 $U$ 上把丛 $E$ 平凡化。于是可以认为，$E\to C$ 是把 $U$ 上的平凡丛与 $C\backslash p$ 上的丛 $E$ 用平凡粘合映射粘合而成。然后，把粘合映射改为 $z^\m$（其中 $\m$ 是对偶群的一个整权，$z$ 是 $p$ 处的局部参数），由此产生一个新丛 $E'$。丛 $E$ 与 $E'$ 在 $U$ 之外可以用同一个联络描述。描述 $E$ 的联络 $a$ 在 $U$ 上为零；描述 $E'$ 的联络 $a'$ 在 $U$ 之外与 $a$ 重合，而在 $U$ 内可以取为某个特定的参考联络，它支撑在 $U'$ 中的一个环形区域上，并正比于 $\m$。

取表示输入波函数的泛函 $\Psi$，在 $a'$ 上求值。这就给出一个新的泛函 $\Psi_{\m}(a)$。关键在于，$\Psi_\m(a)$ 在复化规范变换下并不协变：新丛 $E'$ 依赖于最初对 $E$ 的平凡化的选取。形式上，$\Psi_\m(a)$ 及其泛函导数可以解释为 Kac--Moody 流在存在“谱流算符”（spectral flow operator）$\Sigma_{\m}(0)$ 时的关联函数。Hecke 算符与 Kac--Moody 代数谱流自同构之间的这一关系，先前已从另一虽不同但相关的角度被注意到，见 \cite{Teschner:2010je} 第 8.5 节。

“谱流”一词指的是 Kac--Moody 代数的某个自同构：
\begin{align}
J^\alpha_n &\to J^{\alpha}_{n+(\m, \alpha)} \cr
J^h_n &\to J^h_n - \m \kappa \delta_{n,0}
\end{align}
其中 $J^\alpha$ 是与根 $\alpha$ 相联系的流，$J^h$ 是 Cartan 流。这恰是 $z^\m$ 规范变换对各流的 Fourier 模的作用。

按定义，谱流模（spectral flow module）是真空模在谱流下的像。特别地，真空向量在谱流下的像被 $n$ 非负的 $J^{\alpha}_{n+(\m, \alpha)}$ 消灭，并且是 $J^h_0$ 的一个具有非平凡本征值的本征向量。\zy{译者按（无外部依据）：谱流并不保持真空，它把整条能谱沿 $J^h_0$ 平移；谱流模因此是真空模在这种平移下的像，其最低权态不再是真空，这正是下文 $|1\rangle$ 满足 $J^0_0|1\rangle=|1\rangle$ 的原因。}相应地，谱流算符 $\Sigma_{\m}(0)$ 是这样的局部算符：在作 $z^\m$ 规范变换之后，它与 Kac--Moody 流的 OPE 变得非奇异。$J^{\alpha}(z)$ 在 $z=0$ 处将有 $(\m, \alpha)$ 阶的极点或零点，而 $J^h(z)$ 将有一个单极点，其留数为\footnote{玻色化（bosonization）为描述 $\Sigma_{\m}$ 提供了一种便利的方式。粗略地说，若把 Cartan 流玻色化为 $J^h = \partial \phi^h$，把其余的流写成顶点算子（vertex operator）$J^\alpha = e^{\frac{\alpha}{\kappa} \cdot \varphi}$，则谱流算符也可以用顶点算子 $\Sigma_{\m} = e^{\m \cdot \varphi}$ 表示。这一表示对某些计算有用，但在一般的 $G_\C$ 规范变换下表现不佳。}$\m\kappa \Sigma_{\m}(0)$。

我们再次强调，泛函 $\Psi_{\m}(a)$ 并不表示 $\M(G,C)$ 上的半密度，因为它依赖于对 $E$ 的平凡化的选取。谱流算符 $\Sigma_{\m}(0)$ 的性质，恰好刻画了规范协变约束（\ref{eq:cogauge}）失效的确切方式。我们的目标是由 $\Sigma_{\m}(0)$ 构造某种局部算符插入，它要是中心的，从而能表示 Hecke 算符对 $\Psi(a)$ 的作用。

在继续一般性讨论之前，先给出 $G=\SO(3)$、最小电荷这一参考例子。基本的谱流自同构为
\begin{equation}
J^\pm_n \to J^{\pm}_{n \pm 1} \qquad \qquad J^0_n \to J^0_n - \delta_{n,0}
\end{equation}
谱流模由一个向量 $|1\rangle$ 生成，它满足
\begin{align}\label{laterref}
J^{\pm}_{n \pm 1} |1\rangle &=0  \qquad \qquad n\geq 0 \cr
J^0_n |1\rangle &=0  \qquad \qquad n> 0 \cr
J^0_0|1\rangle &= |1\rangle.
\end{align}
在这种基本的“电荷 1”谱流算符的情形，我们把 $\Sigma_\m(0)$ 记作 $\Sigma_1(0)$。

依照第 \ref{affgr} 节关于仿射 Grassmann 流形的讨论，我们可以把粘合映射 $z^\m$ 换成同一轨道 $\Gr^\m$ 中的另一个粘合映射 $g$。把 $z^\m$ 换成 $g$ 作同样的构造，便产生一个泛函 $\Psi_{g}(a)$。$\Psi_{g}(a)$ 的泛函导数可以解释为在存在修改后的谱流算符 $\Sigma_{g}(0)$ 时的关联函数。由于平凡化的改变由各流的非负模实现，我们可以把这些模对 $\Sigma_{g}(0)$ 的作用表示成沿 $\Gr^\m$ 的某些微分算子。

应当强调，$\Psi_{g}(a)$ 的定义确实需要选取一个支撑在 $U'=U\backslash p$ 内、实现粘合映射 $g$ 的参考联络。描述同一 $g$ 的不同联络之间差一个复规范变换，因而可能给出 $\Psi_{g}(a)$ 与 $\Sigma_{g}(0)$ 的不同归一化。结果是，$\Psi_{g}(a)$ 与 $\Sigma_{g}(0)$ 实际上是 $\Gr^\m$ 上某个线丛的截面。我们将在第 \ref{fft} 节辨认出这个线丛。\zy{译者按（无外部依据，据原文校勘）：原文此句作 “indentify”，当为 “identify”（辨认、认出）。译文按后者处理。}各流的非负模将作为这个线丛截面上的向量场作用。这正是 Hecke 修正前后半密度丛不相匹配这一事实在手征代数中的体现。

$\Gr^\m$ 上的这个线丛由 Kac--Moody 代数的 level 控制。下一节我们将说明，在临界 level，这个线丛与 $\Gr^\m$ 上的密度丛重合。这意味着，在一个全纯与反全纯 Kac--Moody level 都取临界值的理论中，谱流算符是 $\Gr^\m$ 上的一个密度，可以自然地积分：
 \begin{equation}\label{tendef}
\h\Sigma_\m(0) \equiv \int_{\Gr^\m}\Sigma_{g}(0) |\d g|^2.
 \end{equation}
$\h\Sigma_\m(0)$ 具有恰当的性质，可作为电荷为 $\m$ 的 Hecke 算符，它按第 \ref{wilop} 节的说法对偶于带有最小角 $s_{R}$ 的 Wilson 算符（即用最高权向量构造的波函数所定义的 Wilson 算符）。对 $\Gr^\m$ 的积分推广了式 (\ref{tonf}) 中对 $\bCP^1_x$ 的积分。\zy{上文的临界 level（critical level）指 Kac--Moody 的 level 取 $\kappa_c=-h$，其中 $h$ 是对偶 Coxeter 数；此时 Sugawara 构造退化，手征代数的中心变得很大，故有 $S_{\cal P}(z)$ 这样的中心元。参见 B. Feigin 与 E. Frenkel, ``Affine Kac--Moody algebras at the critical level and Gelfand--Dikii algebras,'' \emph{Int. J. Mod. Phys.} A7 (1992)。}
 
各流的非负模对 $\h\Sigma_\m(0)$ 的作用，可以换成 $\Sigma_{g}(0)$ 沿 $\Gr^\m$ 上相应向量场的 Lie 导数。只要分部积分时不出现边界项（在可能出现单极子冒泡时，这或许需要作技术性的分析），$\h\Sigma_\m(0)$ 就会被 Kac--Moody 代数的非负模消灭，从而是中心的。相应地，平均后的泛函
 \begin{equation} \label{eq:chirhecke}
\int_{\Gr^\m}\Psi_{g}(a)|\d g|^2
 \end{equation}
它由 $\h\Sigma_\m(0)$ 作用于 $\Psi(a)$ 而得，是规范协变的，可以表示电荷为 $\m$ 的主 Hecke 算符（principal Hecke operator）的作用。

我们可以把这套构造直接用于前面 $G=\SO(3)$、最小电荷的示例。可以定义一族以 $\bCP^1$ 参数化的谱流算符 $\Sigma_{1;\mu}(0)$，即把 $\Sigma_1(0)$ 作整体 $\SO(3,\C)$ 旋转。形式上，可以把相应的态写成
\be 
|1;\mu\rangle = e^{\mu J^+_0 }|1\rangle
\ee
把非负 Fourier 模对 $|1;\mu\rangle$ 的作用写成关于 $\mu$ 的微分算子，并验证它们都是全导数，这是不难的。显然有
\be 
J^+_0 |1;\mu\rangle = \partial_\mu e^{\mu J^+_0 }|1\rangle=  \partial_\mu |1;\mu\rangle
\ee
$J^0_0$ 的作用同样不难写出
\be 
J^0_0 |1;\mu\rangle = e^{\mu J^+_0 }(J^0_0 + \mu J^+_0)|1\rangle  =  \partial_\mu \left(\mu |1;\mu\rangle\right)
\ee
计算 $J^-_0$ 的作用只需稍多一点工夫：
\be 
J^-_0 |1;\mu\rangle = e^{\mu J^+_0 }(J^-_0 + 2 \mu J^0_0 + \mu^2 J^+_0)|1\rangle  =  \partial_\mu \left(\mu^2 |1;\mu\rangle\right)
\ee

这就使得下述插入
\be \label{eq:averso}
\int \Sigma_{1;\mu, \bar \mu}(0)|\d \mu|^2
\ee
是中心的，只要分部积分时边界项消失。说明边界项消失的自然办法，是说明 (\ref{eq:averso}) 确实是 $\bCP^1$ 上某个密度的积分。为此，我们需要用第二个坐标片覆盖 $\bCP^1$：从相反的谱流算符 $\Sigma_{-1}(0)$ 出发，把它形变为 $\Sigma_{-1;\mu}(0)$，即
\be 
|-1;\mu\rangle = e^{-\mu^{-1} J^-_0 }|-1\rangle
\ee
非负 Fourier 模对这一族算符的作用，涉及的关于 $\mu$ 的微分算子与作用在 $\mu^2 |1;\mu\rangle$ 上的相同。把反手征模也算进来，我们发现可以一致地认定
\begin{equation}
\Sigma_{-1;\mu}(0) = |\mu|^4 \Sigma_{1;\mu}(0)
\end{equation}
并把它们合并为定义在整个 $\bCP^1$ 上的一个密度 $\Sigma_{1;\mu}(0)$。第 \ref{fft} 节将给出通向这一结果的另一种更一般的途径。

按构造，插入 $\Sigma_{1;\mu}(0)$ 对应于背景联络的一个非常具体的修改 $a \to a'[a;\mu]$。回忆我们所处的规范是：$a$ 在开坐标片 $U$ 内为零，而 $a'[a;0]$ 与 $a$ 相差一个支撑在 $U'$ 中环形区域上的参考联络。插入取指数后的 Fourier 模 $e^{\mu J^+_0}$，又在更宽的一个环形区域上对联络添加一个具体的修改，从而得到 $a'[a;\mu]$。

于是，与 (\ref{eq:averso}) 相应的积分算符 (\ref{eq:chirhecke}) 为
 \begin{equation}\label{eq:chirhesutwo}
\int_{\bCP^1}\Psi(a'[a;\mu])|\d \mu|^2
 \end{equation}
把它与 (\ref{actst}) 比较。那里的 $\Psi(y)$，在这里应写作 $\Psi(a(y))$，其中 $a(x)$ 表示我们为每个丛 $x$ 选定的规范固定代表联络。没有理由认为 $a'[a(x);\mu]$ 已经处于规范固定的形式。需要作一个复化规范变换，才能把它化为修改后的丛 $x_\mu$ 的规范固定形式 $a(x_\mu)$。反常将给出某个重标度因子，我们可以把它写成某个全纯量 $\omega$ 的绝对值：
  \begin{equation} \label{eq:omega}
\Psi(a'[a(x);\mu]) = |\omega(x;\mu)|^2 \Psi(a(x_\mu))
 \end{equation} 
该积分算符变为
  \begin{equation}
\int_{\bCP^1}|\omega(x;\mu)|^2 \Psi(a(x_\mu))|\d \mu|^2
 \end{equation}
我们得到：
  \begin{equation}
F(x,y) = \int_{\bCP^1} |\omega(x;\mu)|\delta(y;x_\mu)|\d \mu|^2
 \end{equation}
其中 $\delta(y;x)$ 是支撑在 $\M \times \M$ 中对角线上的 delta 函数。

(\ref{eq:omega}) 的左端是关于 $x$ 的半密度、关于 $\mu$ 的密度，右端则涉及关于 $y=x_\mu$ 的半密度。于是我们可以把 $\omega(x;\mu)$ 与 $(\ref{conf})$ 中的 $w(\vec x;\mu)$ 视为同一对象。因子 $\omega(x;\mu)$ 编码了 $\Psi$ 在复化规范变换下的反常重标度，因而非零。这使我们能把它等同于 \cite{BD} 引入的全纯因子 $k$。\footnote{我们马上就会看到，$\omega(x;\mu)$ 可以在伴随自由费米子的理论中算出。} 我们把 (\ref{eq:chirhesutwo}) 等同于与 $G^\vee$ 的二维表示相联系的 Hecke 算符 $H_{p=0}$，其角取最小者、对应于 $|k|^2$，如式 (\ref{tonf}) 那样。


\subsection{自由费米子技巧}\label{fft}

上一节我们看到，在具有 Kac--Moody 对称性的共形场论中，算符 $\Sigma_g(0)$（其中 $g$ 是与权为 $\m$ 的 Hecke 变换相联系的规范变换）是 $\Gr^\m$ 上某个线丛的截面。这个线丛由反常决定，故只依赖于共形场论的中心荷 $\kappa$（即 Kac--Moody 的 level）。我们希望就临界 level $\kappa_c=-h$ 的共形场论计算这个线丛，但事实证明，对 level 为 $-\kappa_c=+h$ 的共形场论来计算它格外简单。这将给出我们真正想要的 $\Gr^\m$ 上线丛的逆丛。
 
在 $C$ 上选定一个旋量结构——等价地，选定 $K_C$ 的一个平方根 $K_C^{1/2}$——之后，我们考虑一组自旋为 1/2、取值于规范群伴随表示的手征（Majorana--Weyl）费米子 $\psi^a$。Kac--Moody 流是由正规序（normal ordering）的费米子双线性构造而成，并且由于正规序而具有反常的规范变换。中心荷恰为 $h=-\kappa_c$。\zy{正规序指在算符乘积中把产生算符移到消灭算符右侧、并减去真空期望值；它破坏经典对称性，从而产生此处所说的反常规范变换。参见 P. Di Francesco, P. Mathieu and D. Sénéchal, \emph{Conformal Field Theory} (Springer, 1997)。}

我们断言，伴随自由费米子理论中的谱流算符是 $K^{-1}_{\Gr^\m}$ 的截面。这意味着，临界 level $\kappa_c$ 的共形场论中的谱流算符是它的逆丛、即 $K_{\Gr^\m}$ 的截面。我们将以 $G=\SO(3)$、$\m$ 取最小的情形为例加以说明，并简要指出对其他 $G$ 与 $\m$ 的推广。


选定 $\SO(3)$ 的一个 Cartan 子代数后，我们有手征费米子 $\psi^\pm$ 与 $\psi^0$。$z=0$ 处的基本 Hecke 修正使 $\psi^+$ 在 $z=0$ 处有极点、$\psi^-$ 有零点。这只需在 $z=0$ 处插入一个 $\psi^-$ 便可实现。与点 $(u,v)\in\bCP^1$ 相联系的 Hecke 修正，由 $\psi^-$ 的 $\SO(3,\C)$ 旋转实现，即由 $\psi^-_{(u,v)} \equiv u^2 \psi^-  + 2 u v \psi^0 +  v^2 \psi^+$ 实现。插入 $\psi^-_{(y,v)}(0)$ 会迫使 $\psi^-_{(u,v)}$ 在 $z=0$ 处为零，同时使 $\psi$ 各分量的其他线性组合出现极点。\zy{译者按（无外部依据，据原文校勘）：此句下标作 $(y,v)$，与同句及前文的 $(u,v)$ 不一致，当为 $(u,v)$ 之误；译文照录原文。}


于是在伴随费米子的理论中，我们有 $\Sigma_{1;(u,v)}(0) = \psi^-_{(u,v)}$。它关于 $\bCP^1=\Gr^\m$ 的齐次坐标是二次的，因此是 $\O(2) = K_{\Gr^\m}^{-1}$ 的整体截面，与断言相符。于是在临界 level，本例中的谱流算符是 $K_{\Gr^\m}$ 的截面；这一事实已在第 \ref{cvo} 节中用到。

\begin{yizhu}[译者按（无外部依据，据原文梳理）：临界 level 与自由费米子技巧]
本节的论证可概括为一次“取逆”。临界 level 时手征代数的中心很大，谱流算符的归一化只由反常决定，因此它应是 $\Gr^\m$ 上某个线丛的截面，而该线丛只依赖中心荷。直接算 $\kappa_c=-h$ 并不方便，但伴随自由费米子系统的 level 恰为 $-\kappa_c=h$：其流由正规序双线性构造，反常可以显式算出，于是其谱流算符是 $K^{-1}_{\Gr^\m}$ 的截面；取逆便得到临界 level 的答案：线丛 $K_{\Gr^\m}$（全纯典范丛）。这正是把 Hecke 算符写成 $\Gr^\m$ 上积分的原因。
\end{yizhu}


对一般的规范群与电荷，参考 Hecke 修正使以根 $\alpha$ 标记的费米子在 $z=0$ 处多出 $(\lambda, \alpha)$ 阶的极点或零点。这由一个非常简单的顶点算子实现：
\begin{equation}\label{opspin}
\prod_{\alpha | (\lambda, \alpha)<0} \prod_{n_\alpha =0}^{-(\lambda, \alpha)-1}\partial_z^{n_{\alpha}}\psi^{\alpha}
\end{equation}
不难看出，这个乘积按 $K_{\Gr^\m}^{-1}$ 的截面变换：乘积中每个费米子导数，对应一个在 $\Sigma_\m$ 上非平凡地作用的非负 Fourier 模 $J_{n_\alpha}^\alpha$；这些模给出 $\Gr^\m$ 切空间的一组基。

这一计算可以表述为：比较 Hecke 修正前后作用在 $\psi$ 上的 Dirac 算符的 Pfaffian。\cite{BD} 对这一 Pfaffian 作了详细分析。

由于规范反常，算符 $\Sigma_{g}(z,\bar z)$ 并不是 $z$ 的真正函数。事实上，即使把局部坐标重标度 $z \to \lambda z$，也会把奇异规范变换从 $z^\m$ 变为 $(\lambda z)^\m$，从而招致 Cartan 零模 $\m \cdot J_0$ 作用在谱流算符上，产生一个正比于 $(\m,\m)$ 的非平凡标度维数（scaling dimension）。借助自由费米子技巧，可以研究这种对 $z$ 的反常依赖。例如，对 $\SO(3)$ 与最小的 $\m$，我们插入一个 $\psi^-$，它表现为 $K_C^{1/2}$ 的截面。相应地，在临界 level，谱流算符是 $K_C^{-1/2}\otimes \bar K_C^{-1/2}=|K_C|^{-1}$ 的截面。对平均后的 Hecke 算符 $\h\Sigma_\m(0)$ 而言这仍然成立，因为在本例中 $\Gr^\m=\bCP^1$ 上的坐标标度维数为 $0$，故测度 $\d \mu$ 对标度维数没有贡献。$\h\Sigma_\m(0)$ 是 $|K_C|^{-1}$ 的截面这一点，可由 $S$-对偶预期。这与用来定义对偶 Wilson 算符（第 \ref{wilop} 节）的全纯与反全纯截面 $s$ 与 $\bar s$ 分别是 $K_C^{-1/2}$ 与 $\bar K_C^{-1/2}$ 的截面这一事实相符。

对一般的群与表示，把 $\h\Sigma_\m(0)$ 的标度维数与相应 Wilson 算符的行为匹配起来要更微妙一些。费米子插入的标度维数按电荷的二次方增长，而 $\Gr^\m$ 上测度 $\d\mu$ 的负标度维数也按二次方增长。费米子上的导数与测度的标度维数之间有很好的抵消，使得 $\h\Sigma_\m$ 的标度维数是 $\m$ 的一次函数。对偶 Wilson 算符的标度维数已在第 \ref{wilop} 节中描述过。





\subsection{积分–微分 Hecke 算符与 oper 微分方程}\label{opde}

在第 \ref{lineq} 节中，除了积分核以 delta 函数的形式支撑在 Hecke 对应上的积分型 Hecke 算符之外，我们还考虑了积分–微分型 Hecke 算符，其积分核是 delta 函数的导数。构造这类更一般的 Hecke 算符的自然途径，是考虑插入 $\Sigma_{g}$ 的 Kac--Moody 后代（descendant），它们表示在 Hecke 修正所在位置附近的邻域内取泛函导数 $\frac{\delta}{\delta a}$。我们可以只限于由 Kac--Moody 流负模作用所得的后代，因为非负模可以换成关于 $g$ 的导数，而后者在分部积分时会被积掉。

非负模作用在 $\Sigma_{g}$ 的某个后代上，将给出其他后代的 $g$ 导数的某个线性组合。我们需要后代的某种依赖于 $g$ 的组合，使它们按 $\mathrm{Gr}^\m$ 上密度的方式变换，并且使非负 Kac--Moody 模的作用给出关于 $g$ 的全导数。可以给出一个简单的例子：$\partial_z \h\Sigma_\m(z, \bar z)$。事实上，基本谱流算符 $\Sigma_\m(z)$ 的 $z$ 导数恰与 Cartan Kac--Moody 后代重合：
\begin{equation}\label{firstone} \partial_z \Sigma_\m(z)= \frac{1}{2 \pi \i} \oint \frac{\d w}{w-z} \m \cdot J (w) \Sigma_\m(z, \bar z)\equiv  \m \cdot J_{-1} \circ \Sigma_\m(z, \bar z). \end{equation}
为说明这一关系，回忆在关联函数中插入 $\Sigma_\m(z)$ 表示在 $z$ 的邻域 $U$ 内对背景联络作一个具体的修改。随着 $z$ 变动，被修改的联络也变动。这一变动支撑在环形区域 $U'$ 内，可以用一个流的插入（该流与背景联络的变分相缩并）来描述。整个比较都在由 $\m$ 确定的规范群 $\U(1)$ 子群内进行。作为一个小的捷径，我们可以在它们所实现的丛修正这一层面上，比较 $z$ 导数的效果与积分流插入的效果。若 $\Sigma_\m(z)$ 在 $U'$ 上实现规范变换 $g(w) = (w-z)^\m$，则 $z$ 导数 $\partial_z \Sigma_\m(z)$ 实现 $\partial_z  (w-z)^\m = \frac{\m}{z-w} (w-z)^\m$。其中 $\frac{\m}{z-w}$ 部分对应于 $\m \cdot J_{-1}$ Fourier 模所产生的规范变换，而 $(w-z)^\m$ 再次表示 $\Sigma_\m(z)$。由于我们是在丛修正而非联络的层面上工作，可能会漏掉反常引起的效应。在自由费米子理论中作一个简单的检验即可排除这一点。\footnote{对 $\Sigma_\m(z)$ 的玻色化描述，也是验证这一计算的有效办法。}

把这一关系代入 $\h\Sigma_\m(z, \bar z)$ 的定义，我们发现 $\partial_z \h\Sigma_\m(z, \bar z)$ 可以写成 $\Sigma_{g}(z, \bar z)$ 的某个特定 Kac--Moody 后代在 $\Gr^\m$ 上的积分。

产生这类良定义的积分–微分算子的另一种自然途径，是用 Sugawara 向量或其他中心元的模来取 $\h\Sigma_\m$ 的后代。所得的局部算符显然是规范不变的。我们预期，$\Sigma_{g}(z, \bar z)$ 的 Kac--Moody 后代中那些在 $\mathrm{Gr}^\m$ 上为全导数的 $g$ 依赖组合的分类，将与 't Hooft 线端点 $\alpha_{R_\m,n}$ 的相应分类一致。

%This kind of calculation offers a way to prove the oper-like differential equations which we expect to be satisfied by Hecke operators. For example, for $\SO(3)$ and minimal $\m$, which we denote as $\m=1$, we have 
%\begin{align}\label{useform}
%\partial \Sigma_1 &= J_{-1}^0 \Sigma_1 \cr
%\partial^2 \Sigma_1 &= (J_{-2}^0+ J^0_{-1} J^0_{-1}) \Sigma_1 \cr
%S_{-2} \circ \Sigma_1 &= (J_{-2}^0+ J^0_{-1} J^0_{-1} + J^-_{-2} J^+_0) \Sigma_1
%\end{align}
%so that $\partial^2 \Sigma_1 - S_{-2} \circ \Sigma_1$ is proportional to a total derivative $ J^-_{-2} \partial_\mu \Sigma_1$. (The second equation in eqn.
%(\ref{useform}) is just the derivative of the first one, which is eqn. (\ref{firstone}). The last one follows from eqn. (\ref{laterref}) and the definition of $S_{-2}$.)  Thus the averaged spectral flow operator $\h\Sigma_1$ is annihilated by $\partial_z^2 - S_2(z)$, as expected.
%


\subsection{Wakimoto 实现}
这里我们将借助临界 level 的 Wakimoto 实现，给出 $\Sigma_1$ 诸性质的另一种推导。作为额外的收获，我们还将以另一种方式重新得到 oper 微分方程。

临界 level 的 Wakimoto 实现把 Kac--Moody 流表示为 $\bCP^1$ 上扭化 $\beta\gamma$ 系统的对称流：
\begin{align}\label{waki}
J^+ &= \beta \cr
J^0 &= - \beta \gamma + \partial \alpha \cr
J^- &= -\beta \gamma^2 - 2 \partial \gamma + \partial \alpha \gamma \cr
\end{align}
其中 $\alpha$ 是局部定义的全纯函数。Sugawara 向量简化为 Miura 形式 $S_2 = (\partial\alpha)^2 + \partial^2 \alpha$，从而显然是一个恒等算符的倍数，与各流的 OPE 平凡。

谱流自同构自然地延拓到 $\beta\gamma$ 系统，使得谱流算符给 $\gamma$ 一个零点、给 $\beta$ 一个极点。起这种作用的算符通常记作 $\delta(\gamma)$。与由式 (\ref{firstone}) 得到的 $\partial \Sigma_1$ 的预期形式相比较，可得
\begin{equation}
\Sigma_1 = e^{\alpha+ \bar \alpha} \delta(\gamma)\delta(\bar \gamma).
\end{equation}
指数前的因子提供了 $J^0_{-1} \circ \Sigma_1$ 中的 $\partial \alpha$ 部分。

对这一表达式作 $\SO(3,\C)$ 旋转，给出
\begin{equation}
\Sigma_1(z;\mu) = e^{\alpha+ \bar \alpha} \delta(\gamma -\mu)\delta(\bar \gamma- \bar \mu).
\end{equation}
其中 $\mu$ 是 $\bCP^1$ 上的非齐次坐标。对 $\mu$ 的积分很容易算出，结果为
\begin{equation}
\h\Sigma_1(z) = e^{\alpha+ \bar \alpha} .
\end{equation}
这是恒等算符的一个倍数，因而被各流的所有非负模消灭。此外，oper 微分方程显然成立：$\partial_z^2 \h\Sigma_1(z) = S_2(z)\h\Sigma_1(z)$。

我们预期这一模式对所有的 $G$ 与 $\m$ 都成立。临界 Wakimoto 实现给出的中心元，具有由取值于 Cartan 的 $\alpha$ 构造的 Miura oper 的形式。谱流算符将具有 $\beta \gamma$ 系统的谱流算符与 $\alpha$ 的某个函数相结合的形式。平均后的谱流算符将给出 $\beta \gamma$ 系统恒等算符的倍数，再乘以 $\alpha$ 的某些函数，这些函数给出 oper 微分方程解的 Miura 表达式。

